\clearpage
\fancyhead[L]{\small MOLLWEIDE. COMMENTATIONES.}
\hypertarget{native-572-lower}{}
\begin{center}\small Göttingische gelehrte Anzeigen. Stück 71. Seite 705..710. 1814 Mai 2.\\[3pt]\rule{.72\linewidth}{.3pt}\end{center}

\emph{Commentationes mathematico-philologicae tres, sistentes explicationem duorum locorum difficilium, alterius Virgilii, alterius Platonis itemque examinationem duorum mensurarum praeceptorum Columellae. Adjecta est epistola ad v. cl.} J. G. \textsc{Schneider} \emph{de excerptis geometricis Epaphroditi et Vitruvii Rufi scripta ab auctore harum commentationum} \textsc{Carolo Brandano Mollweide}, astron. in acad. Lipsiensi professore. 1813. Leipzig. Bei Cnobloch. 122 Seiten in Octav, nebst einer Kupfertafel.

Je seltener sich gründliche mathematische Kenntnisse und eine innig vertraute Bekanntschaft mit den Schriftstellern und Sprachen des Alterthums in Einer Person vereinigt finden, desto willkommener müssen die Aufklärungen sein, welche der gelehrte Verf. der vorliegenden Aufsätze mit Hülfe der Mathematik über verschiedene dunkle Stellen bei alten Schriftstellern verbreitet. Der erste Aufsatz betrifft die bekannte Stelle in Virgils Gedicht vom Landbau über die Zeiten der zweifachen Honigernte.

\SourceLeaf{573}\begin{quote}\small Bis gravidos cogunt fetus: duo tempora messis \\ Taygete simul os terris ostendit honestum \\ Pleias, et oceani spretos pede repulit amnes, \\ Aut eadem sidus fugiens ubi piscis aquosi \\ Tristior hibernas coelo descendit in undas.

\hfill Georg. IV. 231---235. \end{quote}

Die meisten Ausleger sind darüber einig, dass die Jahrszeit, welche in den zwei letzten Versen bezeichnet wird, der Herbst sei, wo die Plejaden des Morgens untergehen, und dass daher bei dem descendit in undas der \emph{cosmische} Untergang verstanden werden müsse: allein über die Worte sidus fugiens piscis aquosi sind eine Menge verschiedener Erklärungen versucht, ohne dass eine sich allgemeinen Beifall hätte erwerben können, und einige Ausleger haben geradezu gestanden, dass sie so wie sie da stehen unerklärbar sein. Herr \textsc{Mollweide} hatte seine Ansicht von dieser Stelle schon vor geraumer Zeit in der Monatl. Corresp. (1802 Mai) vorgetragen; späterhin lieferte er bei Gelegenheit seiner Ernennung zum Professor der Astronomie in Leipzig in einer besonders gedruckten kleinen Schrift eine Umarbeitung jenes Aufsatzes, welche jetzt mit neuen Vermehrungen in der vorliegenden Sammlung erscheint. Er versteht unter piscis aquosus nicht wie die meisten andern Ausleger die Fische im Thierkreise, sondern den südlichen Fisch, obwohl er einen Hauptgrund, womit er in dem ersten Aufsatz diese Meinung unterstützte (dass nemlich das Beiwort aquosus, welches bei den Fischen im Thierkreise etwas müssig zu stehen scheine, sich auf die Stellung des Sterns erster Grösse im südlichen Fisch am Wasserguss des Wassermanns beziehe), in den spätern Umarbeitungen wieder hat fallen lassen (vermuthlich deswegen, weil sich andere Stellen finden, wo piscis mit demselben Beiworte entschieden den Fisch im Thierkreise bedeutet), und jetzt dieses Beiwort nur als eine Anspielung auf die nasse Jahrszeit nimmt. Das \emph{Fliehen} der Plejaden vor dem südlichen Fisch erklärt er dadurch, dass der cosmische Untergang der Plejaden und der acronychische Aufgang des südlichen Fisches ziemlich nahe zusammenfallen, und zum Beweise, dass das Bild einer Flucht in einem solchen Sinne gebraucht werden könne, bezieht er sich auf eine andere Stelle desselben Gedichts:

\begin{quote}\small Vere fabis satio; tum te quoque, Medica, putres \\ Accipiunt sulci, et milio venit annua cura, \\ Candidus auratis aperit quum cornibus annum \\ Taurus, et adverso \emph{cedens} canis occidit astro.

\hfill Georg. I. 215---218. \end{quote}

\noindent wo cedens dieselbe Bedeutung hat, wie fugiens in der obigen Stelle, und nach Hrn. \textsc{Mollweide} das nahe Zusammentreffen des Abend-Untergangs des Hundes mit dem Früh-Aufgange des Stiers andeutet. Mit so vieler Erudition auch Herr \textsc{Mollweide} seine Erklärung unterstützt, so kann Rec. doch nicht leugnen, dass nach seinem Gefühl das Gezwungene in derselben nicht ganz weggeschafft ist. So gern er diese Erklärung der zweiten Stelle gelten lassen mag und in dem Zusammentreffen des Abenduntergangs eines Sterns mit dem Frühaufgange eines andern, wenn es ein Fliehen des erstern vor dem zweiten genannt wird, \emph{trotz der verschiedenen} Tagszeit ein schönes und wahres poetisches Bild erkennt, so hart scheint es ihm, dasselbe Bild bei dem umgekehrten Fall (der hier in Rede steht) anzuwenden, wo der fliehende Stern des Morgens untergeht und der vertreibende des Abends aufgeht. Im Frühjahr fängt das Sternbild des grossen Hundes, welches bis dahin jeden Abend am westlichen Himmel geglänzt hatte, an, ganz unsichtbar zu werden, ungefähr um die Zeit, wo der Stier, welcher seinerseits eine Zeitlang ganz unsichtbar gewesen war, anfängt, wieder sichtbar zu werden und immer mehr in den Frühstunden den östlichen Himmel zu schmücken: was kann man dagegen haben, wenn ein Dich\SourceLeaf{574}ter dies ein Fliehen des Hundes vor dem Stier nennt? Aber wo bleibt die Wahrheit und Anschaulichkeit des poetischen Bildes, wenn der Dichter den Plejaden eine Flucht vor dem Fomahand beilegt in einer Jahrszeit, wo die Erscheinungen eigentlich auf folgende Art einander folgen: des Abends nach Sonnenuntergang, so bald überhaupt Sterne sichtbar werden, sind beide schon aufgegangen, freilich Fohamond beinahe eine halbe Stunde früher als die Plejaden; allein Fomahand erhebt sich nur wenig über den Horizont, und ist lange vor Mitternacht schon wieder untergegangen, während die Plejaden die ganze Nacht hoch am Himmel glänzen und erst in dem Augenblick unter den Horizont gehen, wo die Sonne aufgeht. So scheint es, wenn obige Stelle als Parallelstelle gebraucht werden darf, würde sie natürlicher die Erklärungsart des Bischofs \textsc{Horsley} rechtfertigen, welcher in den zwei letzten Versen die Gleichzeitigkeit des Abenduntergangs der Plejaden mit dem Früh-Aufgange der Fische (im Frühjahr) findet, wenn nicht dieser Auslegung andere nicht unerhebliche Bedenklichkeiten im Wege ständen. [Die Anzeigen der zweiten und dritten Abhandlung und die Briefe finden sich im IV. Bande dieser Werke unter den Anzeigen von Schriften, die Gegenstände aus dem Gebiete der Geometrie behandeln].
